\section{Reproducibility}
\label{sec:repro}

\paragraph{Evidence chain.}
Each reported number is a \verb|\V{claim-id}| macro generated by the paper data
layer from a canonical bundle in the evidence freeze
\texttt{satsa-evidence-freeze-v2}; the build verifies each bundle's hashes against
the freeze before reading it. Tables and figures are generated the same way and
record their source bundles and manifest hashes. An audit script regenerates paper
data, tables and figures and compares them byte for byte with the committed files,
rejects undefined macros and hand-typed decimals, checks that every citation is a
verified row of the literature matrix, and scans for prohibited claim wording.
The package is archived as a write-once publication snapshot that records the
hash of the freeze it was generated from and the hash of every file.

\paragraph{Reproduction levels.}
We classify what has actually been demonstrated:
\begin{itemize}
\item \emph{Level 1, artifact verification (demonstrated):} the freeze verifier
  recomputes every recorded hash; it has passed in the working repository and in a
  fresh clone, and a mutation drill showed that edited manifests, raw results,
  metrics and configurations are rejected.
\item \emph{Level 2, result regeneration (demonstrated):} paper tables, figures and
  values regenerate byte-identically from the frozen bundles, and the freeze's
  exports regenerate identically.
\item \emph{Level 3, experiment reproduction (demonstrated for deterministic
  experiments):} re-running the integrity mutation matrix (T01) and the peer sweep
  (P01) from the recorded code under new experiment identifiers reproduced their
  analytical outputs; timing-based experiments reproduce only in distribution.
\item \emph{Level 4, environment reproduction (not demonstrated):} a container
  image and compose topology exist but were not executed; all runs used one
  Windows machine with the project's Python dependencies.
\item \emph{Level 5, external-data reproduction (demonstrated once):} the external
  experiment was executed twice from the public dataset (X02, since superseded, and X02b, at different
  commits) with identical analytical outputs; it has not been reproduced by a third
  party.
\end{itemize}

\paragraph{Commands.}
All commands are Python scripts in the repository's \texttt{scripts} directory.
Level~1 runs \texttt{verify\_\allowbreak evidence\_\allowbreak freeze.py} on the
freeze; Level~2 runs \texttt{build\_\allowbreak paper\_\allowbreak assets.py} and then
\texttt{audit\_\allowbreak paper.py}; Level~3 re-runs
the experiment runners with each bundle's recorded configuration. The supplementary
material lists the full commands, seeds, dataset provenance and what cannot be
reproduced.

\paragraph{Code, data and artifact availability.}
The SAT-SA source code, experiment runners, evidence freeze, paper data layer and
the literature-search scripts are in the project repository under the MIT licence
(\url{https://github.com/PrathamKapoor/SAT-SA-with-PQC}); the commit that
contains this manuscript and the publication snapshot is identified in the
supplementary material. All controlled data are synthetic and regenerate from the
recorded seeds and configurations. The external dataset~\cite{amaral2018uci498}
is public under a Creative Commons Attribution (CC BY) licence and is not
redistributed; the experiment downloads it
and verifies its checksum. No personal data, no data from a real SOC and no
human-participant data were used.
